\chapter{Control Layers}\label{chap:layers}

Control layers (SCPs, RCPs, permission boundaries) add additional
deny gates on top of the base policy evaluation.  The key property:
layers can only block, never grant.

\section{Definitions}

\begin{definition}[allowsLayered]\label{def:allowsLayered}
\lean{allowsLayered}
\leanok
Evaluates a policy with SCP, RCP, and permission-boundary layers.
Returns an \texttt{allowed} field (true only if the base \texttt{allows}
is true and no layer blocks) and a \texttt{blocked} list.
\end{definition}

\begin{definition}[noLayers]\label{def:noLayers}
\lean{noLayers}
\leanok
The empty layer set: no SCPs, no RCPs, no boundary.
\end{definition}

\section{Theorems}

\begin{theorem}[evalLayers\_noLayers\_blocked]\label{thm:evalLayers_noLayers_blocked}
\lean{evalLayers_noLayers_blocked}
\leanok
With no layers, the blocked list is empty.
\end{theorem}

\begin{theorem}[evalLayers\_blocked\_noContext\_empty]\label{thm:evalLayers_blocked_noContext_empty}
\uses{thm:evalCond_noContext_tu}
\lean{evalLayers_blocked_noContext_empty}
\leanok
If no layer blocks under a concrete context, no layer blocks under
\texttt{noContext} either.  Relies on the fact that Deny conditions
evaluated as T under \texttt{noContext} imply T under any context
(preventing false blocks), and Allow conditions never evaluate to F
under \texttt{noContext} (preventing false no-allow blocks).
\end{theorem}

\begin{theorem}[layers\_narrow]\label{thm:layers_narrow}
\lean{layers_narrow}
\leanok
If \texttt{allowsLayered} allows a request, then the base
\texttt{allows} also allows it.  Layers never grant access the
base policy did not.
\end{theorem}

\begin{theorem}[layer\_add\_monotone]\label{thm:layer_add_monotone}
\uses{thm:evalLayers_noLayers_blocked}
\lean{layer_add_monotone}
\leanok
If the base policy allows a request, then \texttt{allowsLayered}
with \texttt{noLayers} also allows it.  The empty layer set is
transparent.
\end{theorem}

\begin{theorem}[layered\_nocontext\_conservative]\label{thm:layered_nocontext_conservative}
\uses{thm:layers_narrow, cor:allows_nocontext_conservative, thm:evalLayers_blocked_noContext_empty}
\lean{layered_nocontext_conservative}
\leanok
Layered no-context conservatism: if \texttt{allowsLayered} allows
under any context, it allows under \texttt{noContext}.  Combines
\texttt{layers\_narrow}, \texttt{allows\_nocontext\_conservative},
and \texttt{evalLayers\_blocked\_noContext\_empty}.
\end{theorem}
